\section{Position of the result}
\label{sec:position}

The immediate outcome is a certified extension of a supplied-normal-cut
primitive. ProveIt's original compressed source construction already retained
the exceptional handles, their attaching charts, and the parallelity
interfaces. Its terminal handlebody consumer refused any core containing a
three-handle. A diagnostic script demonstrated promising deletions, but did
not provide independent source replay. We prove the relative geometric fact
needed for those deletions, reconstruct the missing cap incidences from the
source, and compose the resulting proof with the existing terminal
certificates.

This is consequential because a source interior vertex is an artefact of the
triangulation, not evidence of difficult knot topology. A recognizer should
not lose its compressed terminal operation merely because such a vertex was
introduced. The universal source-tree theorem shows that this obstruction
can be removed uniformly, without searching normal-coordinate supports or
expanding repeated sheets.

\subsection{The quasipolynomial objective and its present status}

For an input knot diagram with $n$ crossings, the intended target is a
deterministic bound of the form
\[
 \exp\bigl((\log(n+2))^{O(1)}\bigr),
\]
with $n^{O(\log n)}$ being the particularly strong form announced in 2021.
The announcement and a proved, implemented end-to-end complexity theorem
must be distinguished. Lackenby's July 2026 preprint presents an algorithm
using incompressible surfaces, boundary patterns, and hierarchies. Its
Section~9 explicitly withholds a running-time estimate for insufficiently
specified steps. Proposition~9.1 bounds the number of steps by
$L(g+1)^L$, in the paper's notation; Proposition~10.2 bounds a refined
partial hierarchy length by $4c_q(\mathcal H)$. A linear hierarchy bound is
not a logarithmic depth bound. Polynomial encoding and verification of an
essential hierarchy in Theorem~14.1 do not give a polynomial procedure for
discovering it~\cite{lackenby2026,oxford2021}.

Our mathematical target is therefore a carefully delimited part of that
programme. A bounded compressed core must support inexpensive geometric
operations with enough source information to justify their use after a cut.
The present top-handle operation meets that requirement for the native
parallelity interfaces. It does not supply arbitrary inherited-pattern
transport, a complete useful-surface search, or a global recurrence bounding
all intermediate representations.

\subsection{What is reused and what is contributed}

Removing three-handles from a compact connected three-manifold with nonempty
boundary is classical. Jackson gives the needed unmarked statement and a
dual-path proof in Lemma~3.2 of~\cite{jackson2023}. The same paper explains
why recovering parallelity-bundle topology is weaker than retaining its
inclusion in the ambient manifold. Our contribution is the source-specific
relative construction, its reusable specialization across all normal cuts,
and its independently verified native implementation. We do not claim a new
general existence theorem for three-handle elimination.

The graph-matroid partition theorem is also classical~\cite{edmondsfulkerson}.
Its application here gives a precise abstract optimization problem and
certificates of optimal residual deficiency. The chain-complex constraint in
\cref{sec:chain}, however, shows why abstract successes cannot automatically
be interpreted as geometric performance gains. This limitation is part of
the research result, not an omitted qualification.

Within ProveIt, the source-parallelity construction, the initial uncut
two/three-handle queue, the stopping-set characterization of positive raw
cancellation, and the saturated-bundle projection are prior work. The ten
incoming packages were audited for these overlaps. In particular,
\emph{Stopping Cores} and \emph{Saturated Bundle Cancellation} already
identify the ten three-handle exclusions and ask for a safe adapter.
\emph{Genus Budgets and Native Interfaces} supplies an initial source
reduction, but not the later normal-cut adapter with its attached
interfaces~\cite{repo,incoming-stopping,incoming-saturated,incoming-genus}.
\Cref{sec:source-audit} records the full targeted overlap audit.

The contributions proved or implemented here are:
\begin{enumerate}
\item A cap-incidence formula for exceptional normal-cut two-handles and a
relative top-handle elimination theorem in the authenticated source charts.
\item A single source tree that specializes to valid traces for every
admissible normal coordinate vector, with coordinate-independent trace size.
\item An exact short-row forest criterion and deficiency formula for the
abstract joint deletion problem, together with an independent prototype.
\item A homological rigidity theorem showing that much of that apparent
forest-choice freedom disappears on genuine geometric inputs.
\item An exact exceptional-row quotient for a fixed forest, which retains
arbitrary long rows and gives a finite pattern kernel.
\item A native version-6 certificate adapter, reproducible coverage gains,
and a controlled local scheduling benchmark with its common costs exposed.
\end{enumerate}

\subsection{Notation and input size}

Let $\mathcal T$ be a finite triangulation of a connected compact orientable
three-manifold $M$ with nonempty boundary. The mathematical source arguments
allow several boundary components. The native validator in this release
accepts the narrower class with exactly one torus boundary. It checks
reciprocal face gluings, orientability, connectedness, sphere or disk vertex
links, circle or interval edge links, and absence of edge reversal. It does
not prove that $M$ came from a knot diagram.

A normal vector $x$ has seven nonnegative integer coordinates in every
tetrahedron, satisfies the matching equations, and uses at most one
quadrilateral type per tetrahedron. Its normal surface is denoted $S_x$.
Zero coordinates and the empty surface are allowed. The metric cut
$M\mathbin{\backslash\!\backslash}S_x$ is the completion after splitting
along the local sheets; this convention also handles one-sided components.
No two-sidedness assumption on $S_x$ is needed for the local results.

\begin{center}
\begin{tabular}{@{}ll@{}}
\toprule
Symbol & Meaning\\
\midrule
$t$ & Number of source tetrahedra\\
$b_x$ & $\sum_{i=1}^{7t}\lceil\log_2(x_i+2)\rceil$, total coordinate bits\\
$K$ & Complementary exceptional handle core of the supplied cut\\
$\ports$ & Complete labelled interface rectangles of the parallelity pieces\\
$q,m$ & Numbers of core three-handles and two-handles\\
$E,C$ & Two-handle row labels and surviving one-handle column labels\\
$G$ & Cap graph, with a formal exterior root $\rho$\\
$H$ & Occurrence graph for rows of capped length at most two\\
$r_G(Y)$ & Graphic rank, $|V(G)|-\comp(G[Y])$\\
\bottomrule
\end{tabular}
\end{center}

The source triangulation takes $O(t\log(t+2))$ bits. The full input size for
the local operation is thus $O(t\log(t+2)+b_x)$, apart from already supplied
source certificates. A word-operation estimate on dense source IDs will
always be distinguished from a bit-complexity bound. In particular, an
$O(t)$ trace is not a claim that reading arbitrarily large coordinates costs
$O(t)$ bit operations.
